\section{Discussion: what the common description transfers}\label{sec:p2-discussion}
\subsection{Shared thermodynamics and parameter compensation}
% P:p2-discussion-01
The common equation of state, reference conventions and frozen water/carbonate description can be used with different amine reaction networks while fitting only amine-specific quantities.
What transfers is that thermodynamic description and a conditional ability to fit selected properties; uniform pressure, caloric and phase accuracy does not follow.
MDEA and AMP illustrate how changing the same shared quantities can move unfitted errors in opposite directions before refitting, and how amine-specific compensation can restore an earlier pressure pattern.
A3 shows that recovery is not automatic even within MDEA.
The bicarbonate-related pattern is an observation across these calculations, not proof of a uniquely defective bicarbonate parameter or Born term.

% P:p2-discussion-02
The different parameter freedoms and observations also matter.
PZ needs a more extensive species network, retains a weak cation--bicarbonate interaction and large local energy sensitivities, and predicts false additional liquids.
DEA's four reactive parameters have small conditional errors and no active bounds, yet large unfitted-pressure errors remain.
AMP's carbamate diameter and enthalpy correction reach limits with nearly flat inward profiles.
A diameter at a bound is evidence about this constrained fit, not a measurement of an ion size.
One solute per chemical class, different data objectives and different phase grids do not distinguish an intrinsic class effect from those other changes.

\subsection{Source disagreement, extrapolation and omitted effects}
% P:p2-discussion-03
Temperature-range separation narrows the interpretation without assigning cause.
MDEA and PZ have larger positive errors in high-temperature blocks that extrapolate measured reaction laws, but MDEA Xiao and DEA's low-temperature held-out subsets already show substantial errors inside those law ranges.
AMP's decreasing signed error with temperature occurs while its fitted hydrolysis enthalpy correction is on a limit.
These observations indicate conditions where the thermodynamic description is inaccurate; they do not isolate a law slope from an interaction, a solvent-property extrapolation or source disagreement.

% P:p2-discussion-04
The source comparisons also remain incomplete.
AMP lacks sufficient matched experimental states to separate the Kundu/Hartono difference from concentration-dependent model error, especially at its coldest isotherm.
DEA source values near one loading suggest laboratory disagreement, but those states differ in temperature and concentration and cannot determine which source is correct.
Gas dilution, assumed cell pressures, PZ heavy-water treatment and omitted byproduct chemistry have distinct origins and should not be collapsed into one explanation for all pressure errors.
The wrong water--DEA mixing-enthalpy sign adds an independent thermodynamic limitation that fitted absorption heats do not erase.

\subsection{Scope of the evidence}
% P:p2-discussion-05
The comparisons support statements about the specified equilibrium and property calculations at their measured and searched states.
They combine calibration, fixed-parameter prediction of observations outside the objective, comparison-only properties and numerical phase verification.
Previously accessed held-out results do not become blind validation merely because their residuals are absent from the objective.
Finite phase searches do not establish a continuous phase boundary, and the PZ counterexample shows that the search may miss a resolved split.
The results do not qualify absorber or stripper design, kinetic or transport coefficients, solvent blends, degradation or solid equilibrium.
These limits define the demonstrated transfer; no common percentage threshold or cross-amine ranking is inferred.
